\clearpage\pagestyle{otMain}\setlength{\otBlocEspacement}{2.5mm}\fontsize{11}{15}\selectfont\corrigesChapitre{13}{Droites, distances et cercles}
\begin{blocLivre}{cor-13-1}\begin{corrige}{1}{Nommer les objets}\demarche{Regarder si le tracé s’arrête à un point, à deux points ou pas du tout.}
Objet 1 : la $\res{\text{droite (AB)}}$, illimitée des deux côtés. Objet 2 : le $\res{\text{segment [CD]}}$. Objet 3 : la $\res{\text{demi-droite [EF)}}$, qui commence en E et passe par F.
\label{fig-cor-13-1}\figureGeo[corrige]{dis-1}
\end{corrige}\end{blocLivre}
\begin{blocLivre}{cor-13-2}\begin{corrige}{2}{Lire les codages}\demarche{Associer chaque signe à une égalité.}
Un trait : $\res{AB=AD}$. Deux traits : $\res{CB=CD}$. Le petit carré indique un angle droit en A : les droites (AB) et (AD) sont $\res{\text{perpendiculaires}}$.
\label{fig-cor-13-2}\figureGeo[corrige]{dis-2}
\end{corrige}\end{blocLivre}
\begin{blocLivre}{cor-13-3}\begin{corrige}{3}{Le milieu}\demarche{Le milieu partage le segment en deux longueurs égales.}
$AM=7:2=\res{3{,}5\ \mathrm{cm}}$. Pour N : $AN=3$ cm et $NB=7-3=4$ cm. Comme $AN\ne NB$, N n’est $\res{\text{pas le milieu}}$.
\label{fig-cor-13-3}\figureGeo[corrige]{dis-3}
\end{corrige}\end{blocLivre}
\begin{blocLivre}{cor-13-4}\begin{corrige}{4}{Vocabulaire du cercle}\demarche{Repérer si le segment passe par le centre.}
[OA] : un $\res{\text{rayon}}$. [BC] : un $\res{\text{diamètre}}$, qui est aussi une corde passant par O. [DE] : une $\res{\text{corde}}$.
\label{fig-cor-13-4}\figureGeo[corrige]{dis-4}
\end{corrige}\end{blocLivre}
\begin{blocLivre}{cor-13-5}\begin{corrige}{5}{Sur le cercle ou non ?}\demarche{Comparer chaque distance au rayon.}
M est $\res{\text{sur le cercle}}$. N est $\res{\text{à l’intérieur}}$ : il appartient au disque. P est $\res{\text{à l’extérieur}}$.
\label{fig-cor-13-5}\figureGeo[corrige]{dis-5}
\end{corrige}\end{blocLivre}
\begin{blocLivre}{cor-13-6}\begin{corrige}{6}{Rayon et diamètre}\demarche{Un diamètre vaut deux rayons.}
Rayon de $\mathcal C_1$ : $4{,}6:2=\res{2{,}3\ \mathrm{cm}}$. Diamètre de $\mathcal C_2$ : $2\times1{,}4=\res{2{,}8\ \mathrm{cm}}$. Le plus grand est $\res{\mathcal C_1}$.
\label{fig-cor-13-6}\figureGeo[corrige]{dis-6}
\end{corrige}\end{blocLivre}
\begin{blocLivre}{cor-13-7}\begin{corrige}{7}{À égale distance}\demarche{Chaque condition de distance correspond à un cercle.}
On trace le cercle de centre A de rayon 2 cm et le cercle de centre B de rayon 3 cm. Ils se coupent en $\res2$ points, en vert, symétriques par rapport à la droite (AB).
\label{fig-cor-13-7}\figureGeo[corrige]{dis-7}
\end{corrige}\end{blocLivre}
\begin{blocLivre}{cor-13-8}\begin{corrige}{8}{Un cercle de diamètre donné}\demarche{Le centre d’un cercle de diamètre [AB] est le milieu de [AB].}
M est à 3 cm de A. Le cercle a pour centre M et pour rayon $\res{3\ \mathrm{cm}}$ ; il passe par A et par B (en vert).
\label{fig-cor-13-8}\figureGeo[corrige]{dis-8}
\end{corrige}\end{blocLivre}
\begin{blocLivre}{cor-13-9}\begin{corrige}{9}{La corde la plus longue}\demarche{Un diamètre est la plus longue des cordes.}
La plus longue est $\res{\text{[AB]}}$ : elle passe par O, c’est un diamètre de 5 cm. Aucune corde ne dépasse un diamètre : une corde de 6 cm est $\res{\text{impossible}}$.
\label{fig-cor-13-9}\figureGeo[corrige]{dis-9}
\end{corrige}\end{blocLivre}
\begin{blocLivre}{cor-13-10}\begin{corrige}{10}{Comparer au compas}\demarche{Le compas compare deux longueurs sans les mesurer.}
On prend l’écartement de [CD] au compas et on le reporte sur les autres segments. On obtient : $\res{CD<AB<EF}$ ; en mesurant ensuite : $3{,}4$ cm, $3{,}6$ cm et $3{,}8$ cm.
\label{fig-cor-13-10}\figureGeo[corrige]{dis-10}
\end{corrige}\end{blocLivre}
\begin{blocLivre}{cor-13-11}\begin{corrige}{11}{L’antenne}\demarche{Comparer chaque distance au rayon.}
La zone couverte est le disque de centre A et de rayon 3 km, en vert. Sont couverts $\res{\text{V1, V3 et V4}}$ ; V4 est exactement sur le cercle. V2 n’est pas couvert.
\label{fig-cor-13-11}\figureGeo[corrige]{dis-11}
\end{corrige}\end{blocLivre}
\begin{blocLivre}{cor-13-12}\begin{corrige}{12}{Le trésor}\demarche{Traduire chaque distance par un cercle.}
Le trésor est sur le cercle de centre P de rayon 4 cm et sur le cercle de centre R de rayon 3 cm. Les cercles se coupent en $\res2$ points, $T_1$ et $T_2$ : il faut creuser à l’un de ces deux endroits.
\label{fig-cor-13-12}\figureGeo[corrige]{dis-12}
\end{corrige}\end{blocLivre}
\begin{blocLivre}{cor-13-13}\begin{corrige}{13}{La chèvre}\demarche{La zone accessible est un disque.}
Elle atteint tous les points du disque de centre O et de rayon 5 m. Elle broute $\res{\text{F1 et F3}}$ ; F2, à $5{,}5$ m, est hors de portée.
\label{fig-cor-13-13}\figureGeo[corrige]{dis-13}
\end{corrige}\end{blocLivre}
\begin{blocLivre}{cor-13-14}\begin{corrige}{14}{Deux cercles qui se touchent}\demarche{Traduire « être sur un cercle » par une distance.}
M est sur le cercle de centre A, donc $AM=3$ cm ; il est sur le cercle de centre B, donc $BM=3$ cm. M est sur le segment [AB] et $AM=MB$ : M est le $\res{\text{milieu de [AB]}}$.
\label{fig-cor-13-14}\figureGeo[corrige]{dis-14}
\end{corrige}\end{blocLivre}
\begin{blocLivre}{cor-13-15}\begin{corrige}{15}{Vrai ou faux}\demarche{Vérifier les deux conditions de la définition.}
$\res{\text{Non}}$ : M est à égale distance de A et de B, mais il n’est pas sur le segment [AB]. Le milieu I de [AB] vérifie $IA=IB=2$ cm. Pour être le milieu, il faut aussi appartenir au segment.
\label{fig-cor-13-15}\figureGeo[corrige]{dis-15}
\end{corrige}\end{blocLivre}
